% ============================================
% MEM-06 (ZH)
% 中文翻译 — SPEC-935-R356
% ============================================
\section*{\textbf{MEM-06} --- 政府牲畜圈}
\markboth{MEM-06 --- 政府牲畜圈}{}

\noindent \textit{莫兰布多斯——塞阿拉的集中营}

\noindent \textbf{1915 —— 塞阿拉的集中营}

\vspace{0.5em}\noindent\rule{\textwidth}{0.5pt}\vspace{0.5em}

\noindent 政府管它们叫“牲畜圈”。

\noindent 这名字很准确。我们就是牲口。

\noindent 福塔莱萨郊外的阿拉加迪索（Alagadiço）收容营是一片用带刺铁丝网围起来的区域，大约两百米宽、三百米长。没有顶棚。没有墙。除了踩实了的泥土、太阳和人群积聚的气味，什么都没有。

\noindent 第一天，我数了数人。数到三百多就数不清了。后来又有更多的人来。再后来还有。第一个星期结束时，我们挤得像运输途中的牲口。

\noindent 每天都有士兵押着更多的逃荒者到来。他们来自四面八方：基谢拉莫宾、博阿维亚任、塞纳多尔·蓬佩乌、基沙达、卡宁德。士兵们叫我们“莫兰布多斯”。这个词从他们嘴里说出来，仿佛是一种瘟疫的名字。

\noindent \textit{“又来了一群破衣人。”}

\noindent \textit{“破衣人，滚开。”}

\noindent \textit{“给这些破衣人点吃的，别让他们全死光了。”}

\vspace{0.5em}\noindent\rule{\textwidth}{0.5pt}\vspace{0.5em}

\noindent 食物每天一次。

\noindent 一勺稀汤。士兵们管它叫“政府汤”。是水加面粉，漂着几片叶子。我喝完自己的，还要舔干净葫芦。母亲把她的那份给我。我说不要。她说她已经吃过了。撒谎。我看见咳嗽摇晃着她空荡荡的身体。

\noindent 政府偶尔送来牛。不是给我们吃的——是给士兵吃的。我们只能看着他们烤肉。气味飘过铁丝网，钻进我们的胃里。有人哭。有人捶打铁丝网。士兵们大笑。

\noindent \textit{“让他们闻吧，”其中一个说，“饥饿闻是闻不好的。”}

\noindent 他说得对。饥饿闻不好。哭不好。祈祷也不好。饥饿只能靠食物治好。

\noindent 而没有食物。

\vspace{0.5em}\noindent\rule{\textwidth}{0.5pt}\vspace{0.5em}

\noindent 死亡在第二个星期开始。

\noindent 先是老人。然后是孩子。然后是来之前就已经生病的成年人。

\noindent 士兵们有一套流程。早晨，他们打开铁丝网，走进来，在人群之间穿行，指出那些在夜里死去的人。他们抬着尸体的胳膊和腿——一边一个——抬向乱葬坑。

\noindent 坑在两百米外。我们能看见秃鹫。

\noindent 母亲咳着。

\vspace{0.5em}\noindent\rule{\textwidth}{0.5pt}\vspace{0.5em}

\noindent 一天夜里，我梦见了食物。

\noindent 是个好梦。我在祖母家里——她在我出生前就死了，但梦里她还活着。她在做玉米糕。玉米糕从锅里升起来，升向天空。我张开嘴，它像雨一样落进我嘴里。

\noindent 我饿醒了。

\noindent 一种疼痛的饿。不在胃里。在整个身体里。骨头疼。嘴巴疼。牙齿疼。我看着自己的手，看见皮下的骨头。

\noindent \textit{“娘，”我说，“我想离开这里。”}

\noindent 她没有回答。

\noindent \textit{“娘？”}

\noindent 她侧躺着。身体滚烫。呼吸又短又快。像一只小鸟。

\noindent \textit{“去吧，”她说，“去吧，若阿金。去吧。”}

\noindent \textit{“去哪里？”}

\noindent 她咳起来。咳嗽摇晃着整个身体。咳完之后，她的嘴唇是红的。

\noindent \textit{“离开这里。”}

\noindent \textit{“我不知道该去哪里。”}

\noindent 她闭上眼睛。

\noindent \textit{“你知道的。”}

\noindent 而我知道。

\noindent 我不想知道。但我知道。

\noindent 乱葬坑。铁丝网之外唯一的东西。唯一不需要走路的方向。乱葬坑。那里有肉。

\noindent 那一夜我没有去。

\noindent 第二天夜里我去了。\NE{阿拉加迪索（Alagadiço）收容营在1915年大旱期间位于福塔莱萨附近，这一点有历史记录。若阿金、他的家人以及本段的感官经历均属虚构；营地规模、日常情形和乱葬坑等细节属于叙事重构，并非档案证词。}

\vspace{0.5em}\noindent\rule{\textwidth}{0.5pt}\vspace{0.5em}

\noindent 阿拉加迪索（Alagadiço）收容营建于1915年大旱期间，位于福塔莱萨郊外。历史记录表明，数千名逃荒者曾被关押在那里，但现有资料不足以为本叙事确定精确的收容或死亡人数。

\noindent 我只知道，很多人没有离开。

\noindent 我是活下来的人之一。

\noindent 有时我想：\textit{为什么是我？}

\noindent 于是怪物在我体内笑了。

\noindent 而我记起，我并不是\textit{尽管}有乱葬坑才活下来的。

\noindent 我是\textit{因为}它才活下来的。

\vspace{0.5em}
\noindent\textit{\textbf{↪ 路线:} \rota{MEM-07} \rotasep \rota{MEM-01}}
